\section{AXI GPIO Register Mapping and LED Control}

\subsection{Key Principle}

The \textbf{pin constraint file (XDC) does NOT determine register addresses.}

Register addresses are defined by:
\begin{itemize}
    \item The AXI GPIO IP core hardware design
\end{itemize}

Pin assignments only determine:
\begin{itemize}
    \item Which physical FPGA pin corresponds to which GPIO bit
\end{itemize}

\subsection{AXI GPIO Register Map (Dual Channel Example)}

For a dual-channel AXI GPIO (2 × 32-bit channels), the register layout is:

\begin{center}
\begin{tabular}{|c|c|c|}
\hline
Offset & Register & Controls \\ \hline
0x00 & DATA (Channel 1) & Bits 0--31 \\ \hline
0x04 & TRI (Channel 1)  & Bits 0--31 \\ \hline
0x08 & DATA (Channel 2) & Bits 32--63 \\ \hline
0x0C & TRI (Channel 2)  & Bits 32--63 \\ \hline
\end{tabular}
\end{center}

This layout is defined by the AXI GPIO IP core and is fixed by hardware.

\subsection{Example: 64 LEDs}

If 64 LEDs are used:

\begin{itemize}
    \item Channel 1 controls LEDs 0--31
    \item Channel 2 controls LEDs 32--63
\end{itemize}

Example bare-metal register definitions:

\begin{verbatim}
#define GPIO_BASE 0x41200000

#define GPIO_DATA_CH1 (*(volatile unsigned int*)(GPIO_BASE + 0x00))
#define GPIO_TRI_CH1  (*(volatile unsigned int*)(GPIO_BASE + 0x04))

#define GPIO_DATA_CH2 (*(volatile unsigned int*)(GPIO_BASE + 0x08))
#define GPIO_TRI_CH2  (*(volatile unsigned int*)(GPIO_BASE + 0x0C))
\end{verbatim}

Example usage:

\begin{verbatim}
GPIO_DATA_CH1 = 0xFFFFFFFF;  // LEDs 0–31 ON
GPIO_DATA_CH2 = 0xFFFFFFFF;  // LEDs 32–63 ON
\end{verbatim}

\subsection{How Bits Map to Physical LEDs}

The signal chain is:

\begin{verbatim}
AXI GPIO register bit
        ↓
gpio_io_o[n]
        ↓
leds_o[n]
        ↓
XDC constraint maps to FPGA package pin
\end{verbatim}

Important:

\begin{itemize}
    \item Register offset is defined by the AXI GPIO IP
    \item Bit index is defined by bus wiring order
    \item Physical pin is defined in the XDC file
\end{itemize}

\subsection{Example XDC Constraint}

\begin{verbatim}
set_property -dict {PACKAGE_PIN T22 IOSTANDARD LVCMOS33} [get_ports {leds_o[0]}]
\end{verbatim}

This maps:

\begin{itemize}
    \item \texttt{leds\_o[0]} → FPGA pin T22
\end{itemize}

It does NOT affect the register address.

\subsection{Three Separate Layers}

\begin{enumerate}
    \item Memory Map (Software Layer)  
    Defined by AXI peripheral base address + register offsets.

    \item Signal Index (FPGA Logic Layer)  
    Defined by HDL wiring (e.g., \texttt{gpio\_io\_o[7:0]}).

    \item Physical Pin (Board Layer)  
    Defined in the XDC constraints file.
\end{enumerate}

\subsection{Core Insight}

\begin{itemize}
    \item Memory address selects the peripheral.
    \item Register offset selects the register.
    \item Bit position selects the signal.
    \item XDC selects the physical pin.
\end{itemize}

These layers are connected, but controlled independently during system design.